\hypertarget{sfa}{}
\hypertarget{stochastic-frontier-analysis-full-derivation-and-its-narrow-role-here}{%
\chapter{Stochastic frontier analysis: full derivation and its narrow
role
here}\label{stochastic-frontier-analysis-full-derivation-and-its-narrow-role-here}}

Derive the half-normal composite-error density, conditional
distribution, conditional moments, likelihood and score; understand the
boundary test, unit invariance, identification limits and numerical
implementation.

\hypertarget{from-a-production-frontier-to-a-portfolio-residual-diagnostic}{%
\section{From a production frontier to a portfolio-residual
diagnostic}\label{from-a-production-frontier-to-a-portfolio-residual-diagnostic}}

The classical production-frontier model says observed log output cannot
exceed a stochastic frontier except through symmetric noise:

\begin{equation}
y_{t} = x_{t}{}^{\prime}\beta  + v_{t} - u_{t}, v_{t} \sim N(0,\sigma _{v}^{2}), u_{t} \sim N^{+}(0,\sigma _{u}^{2}), u_{t} \perp  v_{t}.
\end{equation}

N\textsuperscript{+}(0,\ensuremath{\sigma }\textsubscript{u}\textsuperscript{2}) denotes a zero-mean normal
distribution truncated below at zero, equivalently
\textbar N(0,\ensuremath{\sigma }\textsubscript{u}\textsuperscript{2})\textbar. Define the composite
disturbance

\begin{equation}
\varepsilon _{t} = y_{t}-x_{t}{}^{\prime}\beta  = v_{t}-u_{t}.
\end{equation}

The symmetric component v represents ordinary noise; u is a non-negative
shortfall. For a cost frontier the sign is reversed, \ensuremath{\varepsilon }=v+u, because
inefficiency raises cost. Sign orientation must be specified before
estimation.

In this portfolio application, \ensuremath{\varepsilon } is a factor-model residual and u is
merely a candidate one-sided residual component. Calling it managerial
inefficiency would over-interpret the model: omitted factors, crash
exposure, nonlinear payoffs, stale prices and changing volatility can
also create negative skewness.

\hypertarget{component-densities-and-change-of-variables}{%
\section{Component densities and change of
variables}\label{component-densities-and-change-of-variables}}

Let \ensuremath{\phi } and \ensuremath{\Phi } be the standard normal density and cumulative distribution
function:

\begin{equation}
\phi (z)=(2\pi )^{-1/2}e^{-z^{2}/2}, \Phi (z)=\int _{-\infty}^{z}\phi (s)ds.
\end{equation}

The component densities are

\begin{equation}
f_{v}(v)=\sigma _{v}^{-1}\phi (v/\sigma _{v}),
\end{equation}

\begin{equation}
f_{u}(u)=2\sigma _{u}^{-1}\phi (u/\sigma _{u})\cdot 1{u\geq 0}.
\end{equation}

Because \ensuremath{\varepsilon }=v\ensuremath{-}u, we have v=\ensuremath{\varepsilon }+u. The transformation from (v,u) to (\ensuremath{\varepsilon },u) has
Jacobian determinant one. Independence therefore gives the joint density

\begin{equation}
f_{\varepsilon ,u}(\varepsilon ,u)=f_{v}(\varepsilon +u)f_{u}(u), u\geq 0.
\end{equation}

Substitution yields

\begin{equation}
f_{\varepsilon ,u}(\varepsilon ,u)= [1/(\pi \sigma _{v}\sigma _{u})] exp{-\frac{1}{2}[(\varepsilon +u)^{2}/\sigma _{v}^{2}+u^{2}/\sigma _{u}^{2}]}\cdot 1{u\geq 0}.
\end{equation}

\hypertarget{complete-the-square}{%
\section{Complete the square}\label{complete-the-square}}

Introduce total scale \ensuremath{\sigma } and shape ratio \ensuremath{\lambda }:

\begin{equation}
\sigma ^{2}=\sigma _{v}^{2}+\sigma _{u}^{2}, \lambda =\sigma _{u}/\sigma _{v}.
\end{equation}

Define the conditional-location and conditional-scale quantities

\begin{equation}
\mu _{*}=-\varepsilon \sigma _{u}^{2}/\sigma ^{2}, \sigma _{*}^{2}=\sigma _{u}^{2}\sigma _{v}^{2}/\sigma ^{2}.
\end{equation}

Now expand the left side and verify the key identity:

\begin{equation}
(\varepsilon +u)^{2}/\sigma _{v}^{2} + u^{2}/\sigma _{u}^{2} = (u-\mu _{*})^{2}/\sigma _{*}^{2} + \varepsilon ^{2}/\sigma ^{2}.
\end{equation}

To see why, the coefficient on u\textsuperscript{2} is

\begin{equation}
1/\sigma _{v}^{2}+1/\sigma _{u}^{2} = \sigma ^{2}/(\sigma _{u}^{2}\sigma _{v}^{2})=1/\sigma _{*}^{2},
\end{equation}

and the coefficient on u is

\begin{equation}
-2\mu _{*}/\sigma _{*}^{2} = 2\varepsilon /\sigma _{v}^{2}.
\end{equation}

The remaining constant also agrees because
\ensuremath{\mu }\textsubscript{*}\textsuperscript{2}/\ensuremath{\sigma }\textsubscript{*}\textsuperscript{2}+\ensuremath{\varepsilon }\textsuperscript{2}/\ensuremath{\sigma }\textsuperscript{2}=\ensuremath{\varepsilon }\textsuperscript{2}/\ensuremath{\sigma }\textsubscript{v}\textsuperscript{2}.
This single square-completion step produces both the marginal likelihood
and the conditional distribution.

\hypertarget{deriving-the-composite-error-density}{%
\section{Deriving the composite-error
density}\label{deriving-the-composite-error-density}}

Integrate u out of the joint density:

\begin{equation}
f_{\varepsilon}(\varepsilon )=\int _{0}^{\infty}f_{\varepsilon ,u}(\varepsilon ,u)du.
\end{equation}

After the square completion, the u-integral is the upper probability of
a N(\ensuremath{\mu }\textsubscript{*},\ensuremath{\sigma }\textsubscript{*}\textsuperscript{2}) variable:

\begin{equation}
\int _{0}^{\infty} exp[-(u-\mu _{*})^{2}/(2\sigma _{*}^{2})]du = \sigma _{*}\sqrt{2\pi } \Phi (\mu _{*}/\sigma _{*}).
\end{equation}

Since \ensuremath{\sigma }\textsubscript{*}/(\ensuremath{\sigma }\textsubscript{u}\ensuremath{\sigma }\textsubscript{v})=1/\ensuremath{\sigma } and

\begin{equation}
\mu _{*}/\sigma _{*} = -\varepsilon \sigma _{u}/(\sigma \sigma _{v}) = -\lambda \varepsilon /\sigma ,
\end{equation}

the half-normal SFA composite-error density is

\begin{equation}
f_{\varepsilon}(\varepsilon ) = (2/\sigma ) \phi (\varepsilon /\sigma ) \Phi (-\lambda \varepsilon /\sigma ).
\end{equation}

When \ensuremath{\lambda }=0, \ensuremath{\Phi }(0)=1/2 and the density reduces to
\ensuremath{\sigma }\textsuperscript{\ensuremath{-}1}\ensuremath{\phi }(\ensuremath{\varepsilon }/\ensuremath{\sigma }), the ordinary Gaussian error. Positive \ensuremath{\lambda }
tilts mass toward negative residuals.

\hypertarget{deriving-u-conditional-on-the-observed-residual}{%
\section{Deriving u conditional on the observed
residual}\label{deriving-u-conditional-on-the-observed-residual}}

Bayes\textquotesingle{} rule gives f(u\textbar \ensuremath{\varepsilon })=f(\ensuremath{\varepsilon },u)/f(\ensuremath{\varepsilon }).
Cancelling all terms independent of u leaves a normal kernel restricted
to u\ensuremath{\geq }0:

\begin{equation}
f(u|\varepsilon )= \phi [(u-\mu _{*})/\sigma _{*}] / [\sigma _{*}\Phi (\mu _{*}/\sigma _{*})], u\geq 0.
\end{equation}

Thus

\begin{equation}
u|\varepsilon  \sim N(\mu _{*},\sigma _{*}^{2}) \quad \text{truncated below at }0.
\end{equation}

Let z\textsubscript{*}=\ensuremath{\mu }\textsubscript{*}/\ensuremath{\sigma }\textsubscript{*} and define
the inverse Mills ratio r(z)=\ensuremath{\phi }(z)/\ensuremath{\Phi }(z). Standard truncated-normal
integration gives

\begin{equation}
E(u|\varepsilon )=\mu _{*}+\sigma _{*}r(z_{*}),
\end{equation}

\begin{equation}
\operatorname{mode}(u|\varepsilon )=max(0,\mu _{*}),
\end{equation}

\begin{equation}
Var(u|\varepsilon )=\sigma _{*}^{2}[1-z_{*}r(z_{*})-r(z_{*})^{2}].
\end{equation}

The first expression is the Jondrow-Lovell-Materov-Schmidt conditional
mean. A very negative \ensuremath{\varepsilon } makes \ensuremath{\mu }\textsubscript{*} positive, shifting
conditional mass toward a larger one-sided shortfall. A positive \ensuremath{\varepsilon } makes
\ensuremath{\mu }\textsubscript{*} negative, but truncation ensures E(u\textbar \ensuremath{\varepsilon })
remains non-negative.

\hypertarget{deriving-the-conditional-exponential-score}{%
\section{Deriving the conditional exponential
score}\label{deriving-the-conditional-exponential-score}}

In production SFA, exp(\ensuremath{-}u) is often interpreted as technical efficiency
when y is log output. Here returns have arbitrary reporting units, so
first derive the more general moment for c\textgreater0:

\begin{equation}
E[e^{-cu}|\varepsilon ] = [1/\Phi (\mu _{*}/\sigma _{*})] \int _{0}^{\infty} e^{-cu} \phi [(u-\mu _{*})/\sigma _{*}]du/\sigma _{*}.
\end{equation}

Complete a second square in the exponent:

\begin{equation}
-cu-(u-\mu _{*})^{2}/(2\sigma _{*}^{2}) = -[u-(\mu _{*}-c\sigma _{*}^{2})]^{2}/(2\sigma _{*}^{2}) - c\mu _{*} + c^{2}\sigma _{*}^{2}/2.
\end{equation}

The remaining integral is another truncated normal probability, giving

\begin{equation}
E[e^{-cu}|\varepsilon ] = exp(-c\mu _{*}+c^{2}\sigma _{*}^{2}/2) \cdot  \Phi [(\mu _{*}-c\sigma _{*}^{2})/\sigma _{*}] / \Phi (\mu _{*}/\sigma _{*}).
\end{equation}

For classical log-output efficiency set c=1. For a return series, use
c=1/\ensuremath{\sigma } and define the scale-free diagnostic score

\begin{equation}
TE_{scale-free}(\varepsilon )=E[exp(-u/\sigma )|\varepsilon ]
\end{equation}

\begin{equation}
= exp[-\mu _{*}/\sigma  + \frac{1}{2}(\sigma _{*}/\sigma )^{2}] \cdot  \Phi [\mu _{*}/\sigma _{*}-\sigma _{*}/\sigma ] / \Phi (\mu _{*}/\sigma _{*}).
\end{equation}

Why is this invariant? If returns are multiplied by a positive constant
a---for example, decimals to percentages---then \ensuremath{\varepsilon }, \ensuremath{\sigma }, \ensuremath{\mu }\textsubscript{*}
and \ensuremath{\sigma }\textsubscript{*} all multiply by a. Every ratio in the last
expression is unchanged. In contrast, exp(\ensuremath{-}u) changes under a unit
conversion and is unsuitable for ranking return series reported on
different scales.

\hypertarget{a-numerical-conditional-shortfall-example}{%
\section{A numerical conditional-shortfall
example}\label{a-numerical-conditional-shortfall-example}}

Take \ensuremath{\sigma }\textsubscript{v}=0.015, \ensuremath{\sigma }\textsubscript{u}=0.010 and an observed
residual \ensuremath{\varepsilon }=\ensuremath{-}0.020. Then

\begin{equation}
\sigma =0.01803, \lambda =0.6667, \mu _{*}=0.006154, \sigma _{*}=0.008321, z_{*}=0.7396.
\end{equation}

Because r(0.7396)=0.3940,

\begin{equation}
E(u|\varepsilon )=0.006154+0.008321(0.3940)=0.009432.
\end{equation}

The conditional mode is 0.006154 and the scale-free exponential score is
approximately 0.6257. If all returns are multiplied by 100,
recomputation produces exactly the same 0.6257 score up to numerical
rounding.

\hypertarget{log-likelihood-and-analytic-score}{%
\section{Log-likelihood and analytic
score}\label{log-likelihood-and-analytic-score}}

For observation t, define z\textsubscript{t}=\ensuremath{\varepsilon }\textsubscript{t}/\ensuremath{\sigma } and
h\textsubscript{t}=\ensuremath{-}\ensuremath{\lambda }z\textsubscript{t}. The log-likelihood contribution
is

\begin{equation}
\ell _{t}(\beta ,\sigma ,\lambda )=log 2-log \sigma +log \phi (z_{t})+log \Phi (h_{t}).
\end{equation}

The sample log-likelihood is
\ensuremath{\ell }=\ensuremath{\sum }\textsubscript{t=1}\textsuperscript{T}\ensuremath{\ell }\textsubscript{t}. Let
R(h)=\ensuremath{\phi }(h)/\ensuremath{\Phi }(h). Holding \ensuremath{\lambda } fixed, differentiation gives

\begin{equation}
\partial \ell _{t}/\partial \varepsilon _{t} = -[z_{t}+\lambda R(h_{t})]/\sigma .
\end{equation}

Since \ensuremath{\varepsilon }\textsubscript{t}=y\textsubscript{t}\ensuremath{-}x\textsubscript{t}\ensuremath{{}^{\prime}}\ensuremath{\beta } and
\ensuremath{\partial }\ensuremath{\varepsilon }\textsubscript{t}/\ensuremath{\partial }\ensuremath{\beta }=\ensuremath{-}x\textsubscript{t},

\begin{equation}
\partial \ell _{t}/\partial \beta  = x_{t}[z_{t}+\lambda R(h_{t})]/\sigma .
\end{equation}

Holding \ensuremath{\varepsilon } and \ensuremath{\lambda } fixed, \ensuremath{\partial }z/\ensuremath{\partial }\ensuremath{\sigma }=\ensuremath{-}z/\ensuremath{\sigma } and \ensuremath{\partial }h/\ensuremath{\partial }\ensuremath{\sigma }=\ensuremath{\lambda }z/\ensuremath{\sigma }, so

\begin{equation}
\partial \ell _{t}/\partial \sigma  = [-1+z_{t}^{2}+\lambda z_{t}R(h_{t})]/\sigma .
\end{equation}

Finally, holding \ensuremath{\varepsilon } and \ensuremath{\sigma } fixed, \ensuremath{\partial }h/\ensuremath{\partial }\ensuremath{\lambda }=\ensuremath{-}z, hence

\begin{equation}
\partial \ell _{t}/\partial \lambda  = -z_{t}R(h_{t}).
\end{equation}

If optimisation uses \ensuremath{\eta }=log \ensuremath{\sigma } and \ensuremath{\kappa }=log \ensuremath{\lambda }, apply the chain rule:
\ensuremath{\partial }\ensuremath{\ell }/\ensuremath{\partial }\ensuremath{\eta }=\ensuremath{\sigma }\ensuremath{\partial }\ensuremath{\ell }/\ensuremath{\partial }\ensuremath{\sigma } and \ensuremath{\partial }\ensuremath{\ell }/\ensuremath{\partial }\ensuremath{\kappa }=\ensuremath{\lambda }\ensuremath{\partial }\ensuremath{\ell }/\ensuremath{\partial }\ensuremath{\lambda }. A pure log \ensuremath{\lambda } parameterisation cannot
represent \ensuremath{\lambda }=0, so estimate the Gaussian null separately with
\ensuremath{\sigma }\textsubscript{u} fixed at zero.

\hypertarget{numerical-implementation-that-can-survive-difficult-tails}{%
\section{Numerical implementation that can survive difficult
tails}\label{numerical-implementation-that-can-survive-difficult-tails}}

\begin{verbatim}
from scipy.special import log_ndtr
from scipy.stats import norm

def loglik(theta, y, X):
    beta = theta[:-2]
    sigma = np.exp(theta[-2])       # positive total scale
    lam = np.exp(theta[-1])         # positive under alternative
    eps = y - X @ beta
    z = eps / sigma
    return np.sum(np.log(2.0) - np.log(sigma)
                  + norm.logpdf(z) + log_ndtr(-lam * z))

# Fit Gaussian null separately: lambda = sigma_u = 0.
# Check analytic score against central finite differences.
# Use multiple starts and record optimiser convergence/gradient norm.
\end{verbatim}

\begin{itemize}
\tightlist
\item
  Use log \ensuremath{\Phi } via \texttt{log\_ndtr}; directly evaluating \ensuremath{\Phi } in a far tail
  can underflow to zero.
\item
  Evaluate inverse Mills ratios with stable log-density minus log-CDF
  calculations.
\item
  Parameterise positive scales logarithmically and transform back only
  for reporting.
\item
  Check the analytic gradient against a high-accuracy central numerical
  gradient at random interior parameter values.
\item
  Simulate known parameters, refit, and verify bias, coverage, null
  rejection and unit invariance.
\end{itemize}

\hypertarget{boundary-likelihood-ratio-testing}{%
\section{Boundary likelihood-ratio
testing}\label{boundary-likelihood-ratio-testing}}

The null H\textsubscript{0}:\ensuremath{\sigma }\textsubscript{u}=0 is on the boundary
\ensuremath{\sigma }\textsubscript{u}\ensuremath{\geq }0. Under the null, \ensuremath{\lambda }=0 and the model collapses to a
Gaussian disturbance. Let \ensuremath{\ell }\textsubscript{1} be the maximised SFA
log-likelihood and \ensuremath{\ell }\textsubscript{0} the maximised Gaussian-null
log-likelihood:

\begin{equation}
LR = 2(\ell _{1}-\ell _{0}) \geq  0.
\end{equation}

The usual \ensuremath{\chi }\textsuperscript{2}\textsubscript{1} reference assumes an interior null and is
invalid. Under standard boundary regularity conditions,

\begin{equation}
LR \Rightarrow  \frac{1}{2}\chi ^{2}_{0} + \frac{1}{2}\chi ^{2}_{1},
\end{equation}

where \ensuremath{\chi }\textsuperscript{2}\textsubscript{0} is a point mass at zero. Therefore

\begin{equation}
p=1 \text{ if } LR=0; p=\frac{1}{2}\cdot P(\chi ^{2}_{1}\geq LR) \text{ if } LR>0.
\end{equation}

There are 25 portfolio-specific SFA boundary tests, so apply BH FDR to
those 25 p-values as one prespecified family. Only portfolios with
q\ensuremath{\leq }0.05 receive an \texttt{sfa\_rank}; all others are explicitly marked
unsupported rather than quietly ordered.

\hypertarget{identification-why-shape-not-merely-location-matters}{%
\section{Identification: why shape, not merely location,
matters}\label{identification-why-shape-not-merely-location-matters}}

For the half-normal distribution, E(u)=\ensuremath{\sigma }\textsubscript{u}\ensuremath{\sqrt{2/\pi}}, so

\begin{equation}
E(\varepsilon )=E(v-u)=-\sigma _{u}\sqrt{2/\pi }.
\end{equation}

A regression intercept and a non-zero mean shortfall can both explain a
constant downward location shift. If the model contains an unrestricted
intercept, separating that intercept from E(u) relies primarily on
distributional shape---especially skewness---not on the mean alone. That
is weak evidence when residuals are non-normal, serially dependent and
ARCH.

More formally, identification asks whether distinct parameter vectors
generate distinct observable distributions. If a flexible location term,
omitted factors or time-varying volatility can mimic the skewness
attributed to u, the economic label attached to \ensuremath{\sigma }\textsubscript{u} is
not securely identified even when the optimiser returns a precise
number.

\textbf{Interpretation boundary.} A rejected Gaussian boundary null
supports one-sided residual asymmetry under this parametric model. It
does not identify managerial skill, technical efficiency, arbitrage or
mispricing. The factor alpha remains the paper\textquotesingle s primary
performance estimand.

\hypertarget{simulation-verification-plan}{%
\section{Simulation verification
plan}\label{simulation-verification-plan}}

\begin{longtable}[]{@{}p{0.20\textwidth}p{0.34\textwidth}p{0.36\textwidth}@{}}
\toprule\noalign{}
Experiment & Data-generating process & Expected result \\
\midrule\noalign{}
\endhead
\bottomrule\noalign{}
\endlastfoot
Gaussian null & \ensuremath{\sigma }\textsubscript{u}=0 & Boundary-test rejection near its
nominal level after correct mixture calibration \\
Interior SFA & \ensuremath{\sigma }\textsubscript{u}\textgreater0 with known
\ensuremath{\beta },\ensuremath{\sigma }\textsubscript{v},\ensuremath{\sigma }\textsubscript{u} & Parameter bias falls with T;
conditional-score code matches numerical integration \\
Unit rescaling & Multiply y and X\ensuremath{\beta } by 100 & \ensuremath{\beta }/scale transforms
appropriately; scale-free score and ranks stay unchanged \\
Misspecified skewness & Skewed or jump residuals without half-normal u &
SFA may reject, demonstrating that rejection is not unique evidence of
inefficiency \\
ARCH residuals & Gaussian mean with time-varying variance & Assess false
one-sided detections caused by volatility dynamics \\
\end{longtable}

\begin{verbatim}
for replication in range(R):
    v = rng.normal(0, sigma_v, T)
    u = np.abs(rng.normal(0, sigma_u, T))
    y = X @ beta + v - u
    fitted = fit_half_normal_sfa(y, X)
    store(fitted.params, fitted.lr_pvalue, fitted.score)

assert gradient_error < tolerance
assert max_abs(score_decimal - score_percent) < 1e-10
\end{verbatim}

\textbf{Paper result.} Only ME5 BM4 survives the 25-test SFA boundary BH
correction: raw p=0.000254 and q=0.006361. Therefore, the evidence
supports one narrow asymmetry diagnostic, not a general 25-portfolio
efficiency league table.

